\exercisesection{}

\begin{enumerate}
\def\labelenumi{\arabic{enumi})}
\item
  Let E be a finite-dimensional vector space over R, C or H, \(\Phi\) and \(\Phi'\) two nondegenerate positive hermitian forms on E. Assume that \(\mathbf{U}(\Phi) \subset \mathbf{U}(\Phi')\) . Show that \(\Phi\) and \(\Phi'\) are proportional. (Make use of the fact that there exists an orthonormal basis \((e_{1}, \ldots, e_{n})\) for \(\Phi\) that is orthogonal for \(\Phi'\) . The mapping \(u \in \mathcal{L}(\mathrm{E}, \mathrm{E})\) such that \(u(e_{i}) = e_{j}\) , \(u(e_{j}) = e_{i}\) , \(u(e_{k}) = e_{k}\) for \(k \neq i, j\) belongs to \(\mathbf{U}(\Phi)\) .)
\item
  Adopt the notations of Lemma 7. Show that \(\Omega\) has negligible complement in \(\mathbf{GL}(n,\mathbf{K})\) for the Haar measure. (Argue as for Prop. 6.)
\end{enumerate}

¶ 3) Let X be a locally compact space in which a locally compact group H operates on the right, continuously and properly, by \((x,\xi)\mapsto x\xi\) ( \(x\in X\) , \(\xi\in H\) ). Let \(\pi\) be the canonical mapping \(X\to X/H\) . Let \(\rho\) be a continuous representation of H in \(\mathbf{GL}(n,\mathbf{C})\) . Show that for every \(b\in X/H\) , there exist a neighborhood U of b in X/H and a continuous mapping r of \(\overline{\pi}(\mathrm{U})\) into \(\mathbf{GL}(n,\mathbf{C})\) such that \(r(x\xi)=r(x)\rho(\xi)\) for all \(x\in\overline{\pi}(\mathrm{U})\) and \(\xi\in H\) . (One may assume X/H to be compact. Let f be a continuous function \(\geqslant0\) on X, not identically zero on any orbit, and with compact support. Set

\[
r (x) = \int_ {\mathbf {H}} f (x \xi) \rho (\xi) ^ {- 1} d \beta (\xi) \in \mathbf {M} _ {n} (\mathbf {C}),
\]

where \(\beta\) denotes a left Haar measure on H. Show that, for \(f\) and U suitably chosen, \(r(x) \in \mathbf{GL}(n, \mathbf{C})\) for every \(x \in \overline{\pi}^1(\mathrm{U})\) .

\begin{enumerate}
\def\labelenumi{\arabic{enumi})}
\setcounter{enumi}{3}
\tightlist
\item
  Let \(\mathbf{K}\) be a nondiscrete commutative locally compact field. Let \(\mathbf{A}\) be an algebra of finite rank over \(\mathbf{K}\) .
\end{enumerate}

\begin{enumerate}
\def\labelenumi{\alph{enumi})}
\tightlist
\item
  Let \(\mathrm{T}(n, \mathbf{A})\) be the algebra of matrices \((x_{ij}) \in \mathbf{M}_n(\mathbf{A})\) such that \(x_{ij} = 0\) for \(i > j\) . Show that if \(M = (m_{ij}) \in \mathrm{T}(n, \mathbf{A})\) , then
\end{enumerate}

\[
\mathrm{N} _ {\mathrm{T} (n, \mathrm{A}) / \mathrm{K}} (M) = \prod_ {i = 1} ^ {n} \mathrm{N} _ {\mathrm{A} / \mathrm{K}} (m _ {i i} ^ {i - n + 1}).
\]

(Use Lemma 6.)

\begin{enumerate}
\def\labelenumi{\alph{enumi})}
\setcounter{enumi}{1}
\item
  Let \(\mathrm{T}(n,\mathrm{A})^*\) be the group of invertible elements of \(\mathrm{T}(n,\mathrm{A})\) . Let \(M = (m_{ij}) \in \mathrm{T}(n,\mathrm{A})\) . Show that \(M \in \mathrm{T}(n,\mathrm{A})^*\) if and only if \(m_{ii}\) is invertible in \(\mathbf{A}\) for all \(i\) .
\item
  Show that a left Haar measure on \(\mathbf{T}(n,\mathbf{A})^{*}\) is given by
\end{enumerate}

\[
\left(\prod_ {i = 1} ^ {n} \operatorname{mod} \mathrm{N} _ {\mathrm{A} / \mathrm{K}} (m _ {i i}) ^ {i - n - 1}\right) \cdot \bigotimes_ {i \leqslant j} d \alpha (m _ {i j}),
\]

where \(\alpha\) denotes a Haar measure of the additive group of A.

\begin{enumerate}
\def\labelenumi{\alph{enumi})}
\setcounter{enumi}{3}
\tightlist
\item
  Show that if \(\mathbf{N}_{\mathbf{A} / \mathbf{K}} = \mathbf{N}_{\mathbf{A}^0 /\mathbf{K}}\) , then
\end{enumerate}

\[
\Delta_ {\mathrm{T} (n, \mathrm{A}) ^ {*}} \left(\left(m _ {i j}\right)\right) = \prod_ {i = 1} ^ {n} \operatorname{mod} \mathrm{N} _ {\mathrm{A} / \mathrm{K}} \left(m _ {i i}\right) ^ {2 i - n - 1}.
\]

\begin{enumerate}
\def\labelenumi{\arabic{enumi})}
\setcounter{enumi}{4}
\tightlist
\item
  Equip \(\mathbf{R}^{n}\) with the quadratic form \(x_{1}^{2} + x_{2}^{2} + \cdots + x_{n}^{2}\) . Let \(G_{n}\) be the group of displacements of \(\mathbf{R}^{n}\) (Alg., Ch. IX, §6, No.~6, Def. 3) of determinant 1.
\end{enumerate}

\begin{enumerate}
\def\labelenumi{\alph{enumi})}
\item
  Show that \(\mathbf{G}_n\) is unimodular.
\item
  Every element of \(\mathbf{G}_2\) may be uniquely written in the form
\end{enumerate}

\[
(x, y) \mapsto (u + x \cos \omega - y \sin \omega , v + x \sin \omega + y \cos \omega),
\]

where \(u, v\) are in \(\mathbf{R}\) , and \(\omega\) is an element of the group \(\Theta\) of angles of half-lines, isomorphic to \(\mathbf{U}\) . Show that \(du \otimes dv \otimes d\omega\) (where \(d\omega\) is a Haar measure on \(\Theta\) ) is a Haar measure on \(G_2\) .

\begin{enumerate}
\def\labelenumi{\arabic{enumi})}
\setcounter{enumi}{5}
\tightlist
\item
  Let \(\mathbf{P}\) be the set of complex numbers \(z = x + iy\) whose imaginary part is \(>0\) . Show that \(\mathbf{SL}(2,\mathbf{R})\) operates transitively and continuously on the left in \(\mathbf{P}\) by
\end{enumerate}

\[
\left( \begin{array}{c c} a & b \\ c & d \end{array} \right) \cdot z \mapsto \frac {a z + b}{c z + d}
\]

and that P thus becomes a topological homogeneous space for \(\mathbf{SL}(2,\mathbf{R})\) . Show that \(y^{-2}dx dy\) is an invariant measure on P.

¶ 7) Let G be the unimodular group \(\mathbf{SL}(n,\mathbf{R})\) , \(\Gamma\) the subgroup of G formed by the matrices of determinant 1 with integer entries. Show, by induction on n, that every invariant measure \(\mu\) on the homogeneous space \(G/\Gamma\) is bounded and that, for \(f \in \mathcal{K}(\mathbf{R}^{n})\) , one has, with \(\mu\) suitably chosen,

\[
\int_ {\mathbf {R} ^ {n}} f (x) d x = \int_ {\mathrm{G} / \Gamma} \left(\sum_ {z \in \mathbf {Z} ^ {n}, z \neq 0} f (X \cdot z)\right) d \mu (\dot {X})\tag{1}
\]

\((\dot{X} = X\Gamma)\) . (Let \(G'\) be the subgroup of \(G\) leaving invariant the basis vector \((1,0,0,\ldots,0)\) , so that \(G/G'\) may be identified with \(\mathbf{R}^n\) . Let \(G'' = G' \cap \Gamma\) . Let \(H\) be the subgroup of matrices \(\begin{pmatrix} 1 & w \\ 0 & Y \end{pmatrix}\) of \(G'\) for which \(Y \in \mathbf{SL}(n-1, \mathbf{R})\) has integer entries. Then \(H/G''\) is compact and, by the induction hypothesis, \(G'/H\) has a finite invariant measure, therefore by Prop. 12, §2, No.~8, \(G'/G''\) has a finite invariant measure. Applying Exer. 6 of §2, and Exer. 20 b) of A, VII, §4, show that, for \(f \in \mathcal{K}(\mathbf{R}^n)\) ,

\[
a \int_ {\mathbf {R} ^ {n}} f (x) d x = \int_ {\mathrm{G} / \Gamma} \left(\sum f (X \cdot m)\right) d \mu (\dot {X}),
\]

where the summation is over the set of vectors \(m = (m_1, \ldots, m_n)\) whose coordinates are setwise coprime integers (A, VII, §1, No.~2, Th. 1). Applying this to the functions \(f(2x), f(3x), \ldots\) and summing, prove (1). Then, replacing \(f\) by the function \(x \mapsto \varepsilon^n f(\varepsilon x)\) with \(f \geqslant 0\) , and letting \(\varepsilon\) tend to 0, show that every compact subset of \(G/\Gamma\) has measure \(\leqslant 1\) .

\begin{enumerate}
\def\labelenumi{\arabic{enumi})}
\setcounter{enumi}{7}
\tightlist
\item
  \begin{enumerate}
  \def\labelenumii{\alph{enumii})}
  \tightlist
  \item
    Show that there exists on \(\mathbf{S}_{n-1}\) a measure \(\omega_{n-1}\) invariant under \(\mathbf{O}(n,\mathbf{R})\) and, up to a constant factor, only one (regard \(\mathbf{S}_{n-1}\) as a homogeneous space for \(\mathbf{O}(n,\mathbf{R})\) ). b) The mapping \((t,z)\mapsto tz\) permits identifying the topological space \(\mathbf{R}^n -\{0\}\) with the topological space \(\mathbf{R}_+^*\times \mathbf{S}_{n-1}\) . Show that the measure induced on \(\mathbf{R}^n -\{0\}\) by the Lebesgue measure on \(\mathbf{R}^n\) may then be identified, up to a constant factor, with \(t^{n-1}dt\otimes d\omega_{n-1}(z)\) . (Make use of the fact that the Lebesgue measure on \(\mathbf{R}^n\) is invariant under \(\mathbf{O}(n,\mathbf{R})\) and the fact that the Lebesgue measure of the closed ball with center 0 and radius \(r\) is proportional to \(r^n\) .)
  \end{enumerate}
\end{enumerate}

\begin{enumerate}
\def\labelenumi{\alph{enumi})}
\setcounter{enumi}{2}
\item
  Let \(L_{h}\) be the intersection of \(S_{n-1}\) and the hyperplane \(x_{n}=h\) . The non-empty \(L_{h}\) are the intransitivity classes of \(\mathbf{O}(n-1,\mathbf{R})\) in \(S_{n-1}\) . Show that \(\omega_{n-1}\) may be put in the form \(\int\lambda_{h}d\nu(h)\) , where \(\lambda_{h}\) is a measure invariant under \(\mathbf{O}(n-1,\mathbf{R})\) and carried by \(L_{h}\) , and \(\nu\) is a measure on \([-1,1]\) . Identifying \(L_{h}\) with \(S_{n-2}\) by \(z\mapsto he_{n}+\sqrt{1-h^{2}}z\) , one can assume that \(\lambda_{h}=\omega_{n-2}\) . Let \(K_{\theta}\) be the subset of \(S_{n-1}\) defined by \(\sin\theta\leqslant x_{n}\leqslant1\) . By calculating the Lebesgue measure of the set of tz ( \(0\leqslant t\leqslant1\) , \(z\in K_{\theta}\) ), show that \(\omega_{n-1}(K_{\theta})\) is proportional to \(\int_{\theta}^{\pi/2}\cos^{n-2}\varphi d\varphi\) . From this, deduce that if one sets \(h=\sin\theta(-\pi/2\leqslant\theta\leqslant\pi/2)\) , then the measure \(\nu\) may be identified with \(\cos^{n-2}\theta d\theta\) .
\item
  Conclude, by induction on \(n\) , that
\end{enumerate}

\[
d \omega_ {n - 1} = \cos^ {n - 2} \theta_ {1} \cos^ {n - 3} \theta_ {2} \dots \cos \theta_ {n - 2} d \theta_ {1} d \theta_ {2} \dots d \theta_ {n - 1},
\]

where

\[
\begin{array}{r l} & x _ {1} = \sin \theta_ {1} \\ & x _ {2} = \cos \theta_ {1} \sin \theta_ {2} \\ & \dots \\ & x _ {n - 1} = \cos \theta_ {1} \cos \theta_ {2} \dots \cos \theta_ {n - 2} \sin \theta_ {n - 1} \\ & x _ {n} = \cos \theta_ {1} \cos \theta_ {2} \dots \cos \theta_ {n - 2} \cos \theta_ {n - 1} \end{array}
\]

with \(-\pi / 2 \leqslant \theta_{i} \leqslant \pi / 2\) for \(1 \leqslant i \leqslant n - 2, 0 \leqslant \theta_{n-1} < 2\pi\) .

\begin{enumerate}
\def\labelenumi{\arabic{enumi})}
\setcounter{enumi}{8}
\tightlist
\item
  Let \((e_1, e_2, e_3)\) be the canonical basis of \(\mathbf{R}^3\) . Every matrix \(\sigma \in \mathbf{SO}(3, \mathbf{R})\) such that \(\sigma(e_3) \neq e_3\) and \(\sigma(e_3) \neq -e_3\) may be written in a unique way as
\end{enumerate}

\[
\begin{array}{l} \sigma (\varphi , \psi , \theta) = \\ \left( \begin{array}{c c c} \cos \varphi   \cos \psi - \sin \varphi   \sin \psi   \cos \theta & - \cos \varphi   \sin \psi - \sin \varphi   \cos \psi   \cos \theta & \sin \varphi   \sin \theta \\ \sin \varphi   \cos \psi + \cos \varphi   \sin \psi   \cos \theta & - \sin \varphi   \sin \psi + \cos \varphi   \cos \psi   \cos \theta & - \cos \varphi   \sin \theta \\ \sin \psi   \sin \theta & \cos \psi   \sin \theta & \cos \theta \end{array} \right), \end{array}
\]

where \(0 < \theta < \pi\) , \(0 \leqslant \varphi < 2\pi\) , \(0 \leqslant \psi < 2\pi\) (Euler's angles). Show that \(\sin \theta d\theta d\varphi d\psi\) is a Haar measure on \(\mathbf{SO}(3, \mathbf{R})\) . (Identify the homogeneous space \(\mathbf{SO}(3, \mathbf{R}) / \mathbf{SO}(2, \mathbf{R})\) with \(\mathbf{S}_2\) , use Exer. 8, and Th. 3 of §2.)

\begin{enumerate}
\def\labelenumi{\arabic{enumi})}
\setcounter{enumi}{9}
\tightlist
\item
  Let D be the group of displacements of \(\mathbf{R}^{3}\) with determinant 1, the semi-direct product of \(\mathbf{SO}(3,\mathbf{R})\) and the group T of translations. We employ the notation \(\sigma(\varphi,\psi,\theta)\) of Exer. 9, and denote by \(t(\xi,\eta,\zeta)\) the translation with vector \((\xi,\eta,\zeta)\) .
\end{enumerate}

\begin{enumerate}
\def\labelenumi{\alph{enumi})}
\item
  Let H be the closed subgroup of D that leaves \(\mathbf{R}e_{3}\) stable and preserves orientation on \(\mathbf{R}e_{3}\) . Show that H is the product of the group of translations \(t(0,0,\zeta)\) and the group of rotations \(\sigma(\omega,0,0)\) . From this, deduce that the homogeneous space E = D/H, which may be identified with the space of oriented affine lines of \(\mathbf{R}^{3}\) , possesses a measure invariant under D, and, up to a constant factor, only one. (Make use of Exer. 5 a).
\item
  Let \(\mathbf{E}_1\) be the open subspace of \(\mathbf{E}\) formed by the oriented lines not orthogonal to \(e_3\) ; such an oriented line can be determined by the coordinates \(\xi', \eta'\) of its point of intersection with \(\mathbf{R}e_{1} + \mathbf{R}e_{2}\) , and the coordinates \((\sin \varphi \sin \theta, -\cos \varphi \sin \theta, \cos \theta)\) of a direction vector. Show that a measure on \(E_{1}\) invariant under D is given by
\end{enumerate}

\[
\sin \theta | \cos \theta | d \xi^ {\prime} d \eta^ {\prime} d \varphi d \theta .
\]

\begin{enumerate}
\def\labelenumi{\arabic{enumi})}
\setcounter{enumi}{10}
\tightlist
\item
  Let \(G\) be the compact group \(\mathbf{O}(n, \mathbf{R})\) , \(\lambda\) the normalized Haar measure on \(G\) . For \(0 < k < n\) , let \(H\) be the closed subgroup of \(G\) leaving invariant (globally) the subspace \(\mathbf{R}^k\) of \(\mathbf{R}^n\) ; the homogeneous space \(G / H\) , equipped with its quotient topology, may be identified canonically with the Grassmannian \(E = G_{n-1, k-1}(\mathbf{R})\) (GT, VI, §3, No.~6) of the \(k\) -dimensional linear subspaces of \(\mathbf{R}^n\) . There is a measure \(\mu\) on \(E\) invariant under \(G\) , determined up to a constant. For every subspace \(P \in E\) , let \(\sigma_P\) be the image of the Lebesgue measure of \(\mathbf{R}^k\) under an \(s \in G\) such that \(s \cdot \mathbf{R}^k = P\) (independent of the element \(s \in G\) satisfying this relation). Show that \(\mu\) can be chosen in such a way that the following property is satisfied: for every continuous function \(f\) on \(\mathbf{R}^n\) with compact support, and for every \(P \in E\) , one sets \(F(P) = \int_{P} f(x) d\sigma(x)\) ; then
\end{enumerate}

\[
\int_ {\mathrm{E}} \mathrm{F} (\mathrm{P}) d \mu (\mathrm{P}) = \int_ {\mathbf {R} ^ {n}} \| x \| ^ {k - n} f (x) d x.
\]

(If \(\mathbf{P}_0 = \mathbf{R}^k\) , note that one can write \(\int_{\mathbf{E}} \mathrm{F}(\mathrm{P}) d\mu(\mathrm{P}) = c \int_{\mathbf{G}} \mathrm{F}(s \cdot \mathrm{P}_0) d\lambda(s)\) for a suitable constant \(c\) , and on the other hand,

\[
\int_ {\mathrm{G}} f (s \cdot x) d \lambda (s) = c ^ {\prime} \int_ {\mathbf {S} _ {n - 1}} f (\| x \| z) d \omega_ {n - 1} (z)
\]

with the notations of Exercise 8; finally, make use of Exer. 8 b).)

¶ 12) Let K be a commutative field, E a vector space over K, F a linear subspace of E, p the canonical homomorphism of E onto E/F, and A the set of homomorphisms \(f: E/F \to E\) such that \(p \circ f\) is the identity homomorphism of E/F.

\begin{enumerate}
\def\labelenumi{\alph{enumi})}
\item
  If \(f \in A\) and \(h \in \operatorname{Hom}(\mathrm{E} / \mathrm{F}, \mathrm{F})\) , then \(f + h \in A\) . Show that if \(f, f' \in A\) , there exists one and only one \(h \in \operatorname{Hom}(\mathrm{E} / \mathrm{F}, \mathrm{F})\) such that \(f + h = f'\) . One can therefore regard \(A\) as an affine space of which \(\operatorname{Hom}(\mathrm{E} / \mathrm{F}, \mathrm{F})\) is the space of translations.
\item
  Let B be the set of linear subspaces of E that are supplementary to F in E. Show that \(f \mapsto f(\mathrm{E}/\mathrm{F})\) is a bijection of A onto B. One can therefore regard B as an affine space of which \(\operatorname{Hom}(\mathrm{E}/\mathrm{F}, \mathrm{F})\) is the space of translations.
\item
  Show that if \(u\) is an automorphism of \(E\) such that \(u(F) = F\) , then the mapping \(f \mapsto u \circ f\) is an affine bijection of \(A\) onto \(A\) . (Choose an origin in \(A\) , and observe that the mapping \(h \mapsto u \circ h\) is an automorphism of the vector space \(\operatorname{Hom}(E/F, F)\) .) From this, deduce that the mapping \(F' \mapsto u(F')\) is an affine bijection of \(B\) onto \(B\) .
\item
  Assume that K = R and \(\dim E < +\infty\) . Let G be a compact group, and \(\rho\) a continuous linear representation of G in E such that \(\rho(s)(F) = F\) for all \(s \in G\) . Show that there exists an \(F' \in B\) such that \(\rho(s)(F') = F'\) for all \(s \in G\) . (Make use of c) and Lemma 2 of No.~2.) Obtain this result anew using Prop. 1 of No.~1.
\end{enumerate}

